%! TEX root = thesis.tex

\chapter{Introduction}
\label{sec:introduction}

The energy consumption of large language model (LLM) training has grown
dramatically. Training a single frontier model requires millions of GPU-hours
across thousands of accelerators: DeepSeek~V3 consumes 2.79M GPU-hours on 2048
H800 GPUs~\cite{deepseek-ai_deepseek-v3_2024} and Kimi~K2 trains on 15.5T
tokens with a minimum 256-GPU model-parallel group~\cite{kimi_team_kimi_2025}.
At a per-GPU power draw of 700~W and a power usage effectiveness of
1.27~\cite{jegham_hungry_2025}, a single run can emit hundreds of tonnes of
CO$_2$ depending on the local grid mix. As LLM deployment scales, managing
this environmental cost becomes essential for sustainable AI infrastructure.

Carbon-aware temporal workload shifting pauses training when grid carbon
intensity exceeds a threshold and resumes when it falls below a second
threshold. This approach allows a training run to consume cleaner energy
without hardware modification, but the threshold pair $(\theta_p, \theta_r)$
and start date $t_s$ jointly determine the savings-overhead trade-off. Prior
work on temporal shifting for cloud batch jobs reports 3--34\% operational
carbon savings~\cite{wiesner_lets_2021}, and geo-distributed training during
renewable curtailment windows achieves 88--95\% savings across multiple
sites~\cite{wiesner_distributed_2026}. However, none of these works provides a
systematic method for optimizing the hysteresis policy parameters of a
single-site LLM training run. Ad-hoc thresholds predominate, and the value of
the hysteresis margin (the gap between pause and resume thresholds) remains
unexplored.

We address this gap with a discrete-event simulation engine and an adaptive
grid-search optimizer that together search the three-dimensional policy space
$(\theta_p, \theta_r, t_s)$ and return Pareto-optimal configurations. The
optimizer converges within five to six iterations, requiring only
approximately 3,500 simulations per full optimization run (under one minute on
a single CPU core). We evaluate across two state-of-the-art models and five
electricity grid regions using 5-minute-resolution 2025 carbon intensity data
from Electricity Maps. Across tens of thousands of simulation runs, we
establish systematic Pareto frontiers for carbon-aware LLM training with
hysteresis-based pausing.

The results reveal a clear and actionable pattern: the optimal hysteresis
margin is consistently narrow ($\leq 16$~gCO$_2$eq/kWh), and wide hysteresis
degrades the savings-to-overhead ratio in every region and model
configuration. In high-variability grids (Germany, Italy), savings reach
43.4\% and 32.7\% at 174\% and 199\% time overhead. In low-variance grids
(China, United States), savings remain below 10\% regardless of the chosen
policy.

All simulation and optimization code is implemented as a fully client-side
browser application built with SolidJS and TypeScript, with no backend
dependency. The application provides an interactive live simulator, a batch
run dashboard, and an adaptive optimization page with 3D visualization of the
policy space. The complete project is available as open source.

This report is structured as follows. Chapter~\ref{sec:fundamentals} introduces
the necessary background on LLM training, grid carbon intensity, and hysteresis
policies. Chapter~\ref{sec:related_work} surveys related work in carbon-aware
computing. Chapter~\ref{sec:methodology} describes the simulation engine,
policy formulation, scoring, and optimization algorithm.
Chapter~\ref{sec:implementation} presents the software architecture and
implementation. Chapter~\ref{sec:experimental_setup} details the experimental
configuration. Chapter~\ref{sec:results} presents the full results.
Chapter~\ref{sec:validation_verification} describes validation and verification
procedures. Chapter~\ref{sec:conclusion} concludes with a summary of findings
and future work.
